%
In this section, we introduce a novel framework that links the self-attention matrices $\bm{W}_{qk}$ to bilinear forms, enabling us to analyze how an objective function influences their structure.
%
This approach reveals fundamental patterns in $\bm{W}_{qk}$  that are not apparent when examining $\bm{W}_q$ and $\bm{W}_k$ separately.  
%
For example, we prove that autoregressive and bidirectional training induce directional and symmetric structures in the $\bm{W}_{qk}$ matrices.
%
In the following sections, we define and formalize these concepts.
%






\subsection{Interpreting self-attention with bilinear forms}
\label{sec:self-attention-bilinear-form}
%
Self-attention \citep{vaswaniAttentionAllYou2017,radfordLanguageModelsAre2019} is 
a type of score function $A: \mathbf{R}^{N,d} \times \mathbf{R}^{N,d} \rightarrow \mathbf{R}^{N,N}$ that maps a sequence of $N$ token embeddings with dimension $d$ into a matrix of attention scores.
%
Except for the row-wise softmax function $\sigma(\cdot)$, self-attention is a linear transformation of the embedded tokens.
%
In particular, 
%
\begin{equation}
\label{eq:results:self-attention-bilinear-form}
\begin{split}
    \bm{A}(\bm{X}) 
    & = \sigma\left(\frac{1}{\sqrt{d}}\,\hat{\bm{A}}(\bm{X})\right) \\
    & = \sigma\left(\frac{1}{\sqrt{d}}\,\bm{Q} \bm{K}^T\right) \\
    & = \sigma\left(\frac{1}{\sqrt{d}}\,\bm{X} \bm{W}_{qk}\bm{X}^T\right) \,,
\end{split}
\end{equation}
%s
where $\hat{\bm{A}}(\bm{X})$ is the linear part of self-attention (raw unscaled attention scores), $\bm{X} = [\bm{x}^\top_1, \dots, \bm{x}^\top_N] \in \mathbb{R}^{N,d} $ is the sequence of $N$ token embeddings $\bm{x}_i \in \mathbb{R}^d$,  and $\bm{W}_{qk} = \bm{W}_q\bm{W}_k^\top \in \mathbb{R}^{d,d}$.
%
This equation shows that the linear transformation $\bm{W}_q$ and $\bm{W}_k$ are always combined to compute attention scores with one single matrix $\bm{W}_{qk}$.
%
While the matrices $\bm{W}_q$ and $\bm{W}_k$ are defined separately for computational efficiency, this formulation remains mathematically equivalent \citep[see also][]{elhageMathematicalFrameworkTransformer2021, olssonIncontextLearningInduction2022,darAnalyzingTransformersEmbedding2023}.
%

We observe that $\bm{X} \bm{W}_{qk}\bm{X}^\top$ corresponds to a bilinear form (see Definition \ref{def-bilinear-form}). 
%
Specifically, the entry $\hat{\alpha}_{ij} = [\hat{\bm{A}}]_{ij}$ can be formulated in two equivalent ways: (1) as the canonical dot product between a query $\bm{q}_i$ and a key $\bm{k}_j$  (like in standard Transformer models), or (2) as the dot product between tokens $\bm{x}_i$ and $\bm{x}_j$ under the bilinear form $\bm{W}_{qk}$,
%
\begin{equation}
    \hat{\alpha}_{ij} = \langle \bm{q}_i, \,\bm{k}_j \rangle = \langle \bm{x}_i, \bm{W}_{qk} \bm{x}_j \rangle = \langle \bm{x}_i, \bm{x}_j \rangle_{\bm{W}_{qk}} \,.
\end{equation}
%
Intuitively, this equivalence shows that the matrices $\bm{W}_q$ and $\bm{W}_k$ together define an alternative metric in the embedding space, which quantifies the score of $\bm{x}_i$ and $\bm{x}_j$ without explicitly constructing the query and key vectors.
%
Let us consider a specific input token $\bm{x}_i \in \mathbb{R}^d$. 
%
The self-attention layer maps it to an updated token $\hat{\bm{x}}_i \in \mathbb{R}^d$ as follows,
%
\begin{equation}
    \hat{\bm{x}}^\top_i = \bm{x}^\top_i + \sum_{j= 1}^N  \hat{\alpha}_{ij}\,\bm{x}_j\bm{W}_v \,,
\label{eq-prop-self-attention-bilinear-form}
\end{equation}
%
where the coefficients $\{\hat{\alpha}_{ij}\}$ in the second term represent the projection of $\bm{x}_i$ onto $\text{span}\{\bm{X}\}$ (the subspace spanned by $\bm{X} = \{\bm{x}_0^\top, \bm{x}_1^\top, \dots, \bm{x}_N^\top\})$ in the transformed embedding space defined by $\bm{W}_{qk}$,
%
\begin{equation}
    \sum_j\hat{\alpha}_{ij}\, \bm{x}_j =  \sum_j \langle \bm{x}_i, \bm{x}_j \rangle_{\bm{W}_{qk}} \,\bm{x}_j \,.
\label{eq-proof-self-attention-bilinear-form}
\end{equation}
%
Because the softmax function preserves the order of its input values, the ranking of the raw attention scores $\hat{\alpha}_{ij}$ is maintained in the final normalized scores $\alpha_{ij}$,
%
\begin{equation}
     \alpha_{ij} < \alpha_{ij'} \,  \Leftrightarrow \,\hat{\alpha}_{ij} < \hat{\alpha}_{ij'}\quad \forall i,j, j' \,.
\end{equation}
%
As a result, the attention weights ${\alpha_{ij}}$ define a convex combination of the token vectors ${\bm{x}_j}$, meaning the output lies within the convex hull of the input sequence $\text{Conv}(\bm{X}) \subset \text{span}\{\bm{X}\}$ (see also Appendix \ref{supp-math-self-attention-bilinear-form}).
%
We note that $\bm{W}_v$ is a linear transformation that is applied independently to each projection in the sum and thus does not influence our derivation.

% It follows that $\bm{W}_{qk}$ defines the geometric structure that encodes the relevance of the token embedding  $\bm{x}_j$ in predicting $\bm{x}_i$ through projections within the subspace (see Equation \eqref{eq-prop-self-attention-bilinear-form}). 
% %
% Therefore, token prediction depends on learning an optimal bilinear form $\bm{W}_{qk}$ to project tokens in the embedding space effectively.

In the following sections, we demonstrate how this equivalent formulation of self-attention provides a useful framework for analyzing the training of Transformer models.
%
Specifically, we show that the choice of the objective function such as autoregressive prediction \citep{radfordLanguageModelsAre2019} or bidirectional training \citep{devlinBERTPretrainingDeep2019} produces distinct structural patterns in $\bm{W}_{qk}$.
%






\subsection{Deriving the gradients of self-attention with bilinear forms}
\label{sec:self-attention-derivation-gradient}
%
\input{ICML2025/figures/figure-illustration-one-column}
%%
%% 
%
To show the connection between the objective function and the structures of self-attention matrices, we derive a convenient formulation for the weight update of $\bm{W}_{qk}$.
%
We first formulate a sequence modeling problem with self-supervised training as follows. 
%
Let \( U = \{t_1, \dots, t_N\} \) be a sequence of \( N \) tokens. 
%
For each position \( i \), let \( \mathcal{D}_i \subseteq \{1, 2, \dots, N\} \) denote the \emph{conditioning set}, that is, the set of indices corresponding to the tokens used to predict the token \( t_i \).
%
This formulation makes it possible to isolate and analyze the contribution of each token \( t_j \in \mathcal{D}_i \) to the prediction of \( t_i \).
%
Let $\mathcal{L}(U)$ be the negative log-likelihood of each token $t_i$, expressed as,
%
\begin{equation}
    \mathcal{L}(U;\mathcal{W}) = \sum_{i=1}^N \mathcal{L}(t_i)=  - \sum_{i=1}^N \log p_{\mathcal{W}}(t_i \,|\, t_{\mathcal{D}_i}) \,,
\end{equation}
%
where $\mathcal{W}$ is the set of trainable parameters, and $p_{\mathcal{W}}(t_i \,|\, t_{\mathcal{D}_i}) $ are the conditional densities parameterized by the model (see also Appendix \ref{supp-math-self-attention-derivation-gradients}).
%

Here, we demonstrate that the updates to the matrix $\bm{W}_{qk}$ follow a structured pattern: the contribution of the token $t_j $ in the prediction of the embedding $t_i$ results in adding a rank-1 matrix $\bm{K}_{ij}$ to the matrix $\bm{W}_{qk}$.
%
As a result, the total weight update to $\bm{W}_{qk}$ is expressed as a linear combination of these rank-1 matrices.
%
We formalize this observation in the following proposition.
%
\\
\begin{proposition}
\label{prop-gradients-self-attention}
%
(\textbf{The implicit weight update as sum of rank-1 matrices}).
%
Let \( U = \{t_1, \dots, t_N\} \) be a sequence of \( N \) tokens, and let \( \mathcal{L}(U; \mathcal{W}) \) denote the negative log-likelihood of the sequence under a Transformer model with learnable parameters \( \mathcal{W} \),
%
\begin{equation}
    \mathcal{L}(U; \mathcal{W}) = \sum_i \mathcal{L}(t_i) = - \sum_i \log p_{\mathcal{W}}(t_i \mid t_{\mathcal{D}_i})\,,
\end{equation}
%
where \( \mathcal{D}_i \subseteq \{1, \dots, N\} \) is the set of indices used to condition the prediction of token \( t_i \).
%
Let the self-attention function be defined as in Equation~\eqref{eq:results:self-attention-bilinear-form}.
%
Following the gradient of \( \mathcal{L}(U; \mathcal{W}) \), the implicit gradient-based update to the bilinear weight matrix \( \bm{W}_{qk}^l \) at the \( l \)-th self-attention layer is proportionally equivalent to:

\begin{enumerate}[noitemsep]
    \item 
    The sum of contributions from each token \( t_j: j  \in \mathcal{C}_i \) used to predict \( t_i \),
    \begin{equation}
    \label{eq:weight-update-bilinear-form}
    \Delta \bm{W}_{qk}^l \propto \sum_i \sum_{j \in \mathcal{C}_i} \Delta \bm{W}_{qk}^l\Big|_{t_i \leftarrow t_j} = \sum_i \sum_{j \in \mathcal{C}_i} \beta^l_{ij} \bm{K}^{l-1}_{ij}\,,
    \end{equation}
    where \( \mathcal{C}_i \) denotes the set of context indices for predicting \( t_i \).
%\mathcal{L}_3|_{t_i \leftarrow t_j} \,\, \alpha{3N}
%
\[
\Delta \mathbf{W}_{qk} \propto \sum_{i \in \mathcal{P}_j} \sum_{j} \beta_{ij} \mathbf{K}_{ij}
\]
    \item
    The sum of contributions from each prediction target \( t_i: i \in \mathcal{P}_j \) when predicted with \( t_j \),
    \begin{equation}
    \label{eq:weight-update-bilinear-form-prediction}
    \Delta \bm{W}_{qk}^l \propto \sum_{i \in \mathcal{P}_j} \sum_j \Delta \bm{W}_{qk}^l\Big|_{t_i \leftarrow t_j} = \sum_{i \in \mathcal{P}_j} \sum_j \beta^l_{ij} \bm{K}^{l-1}_{ij}\,,
    \end{equation}
    where \( \mathcal{P}_j \) denotes the set of tokens for which \( t_j \) is included in the context.
\end{enumerate}
%
Here, \( \beta^l_{ij} \) is a scalar that quantifies the contribution of the token embedding \( \bm{x}_j \) to the prediction error for \( \bm{x}_i \) at the \( l \)-th layer, and the matrix \( \bm{K}^{l-1}_{ij} \in \mathbb{M}_d \) is a rank-1 matrix defined by the outer product of the token embeddings at the previous layer,
%
\begin{equation}
    \bm{K}^{l-1}_{ij} = \bm{x}^{l-1}_i {\bm{x}^{l-1}_j}^\top\,.
\end{equation}
%
\end{proposition}
%

We provide proof for this proposition with related remarks in Appendix \ref{supp-math-gradients-self-attention}, and an illustrative description of the forward and backward pass in Figure \ref{fig-illustration}.
%


\subsection{How context and prediction impact the gradient differently}
%
Next, we show how the formulation of $\Delta \bm{W}_{qk}$ enables the analysis of the contribution of any given token $t^*$ to the weight updates and how this affects the properties of $\bm{W}_{qk}$.
%
Indeed, Proposition~\ref{prop-gradients-self-attention} indicates that a token $t^*$ impacts the updates of $\bm{W}_{qk}$ differently when serving as context for predicting other tokens or being itself predicted.
%

When a token $t^*$ serves as context ($t_j = t^*$), the embeddings of all predicted tokens contribute to the column space of $\bm{W}_{qk}$, where the update of every $k$-th column $\bm{w}_{\cdot, k}$ is given by
%
\begin{equation}
    \Delta \bm{w}_{\cdot, k} \Big|_{t_j = t^*}= [\bm{x}_{t^*}]_k \left(\sum_{i \in \mathcal{P}_{t}}\beta_{i*}\bm{x}_i\right)\,.
\end{equation}
%
Only the embedding of $t^*$ is instead added to the row space, where the update of every $m$-th row $\Delta \bm{w}_{m, \cdot}$ is given by
%
\begin{equation}
    \Delta \bm{w}_{m, \cdot} \Big|_{t_i = t^*} = \left(\sum_{i \in \mathcal{P}_{t}}\beta_{i*}[\bm{x}_{i}]_k\right)\bm{x}_{t^*}\,.
\end{equation}
%
Intuitively, using $t^*$ as context increases the dimensionality of the column space proportionally to the embeddings of the predicted tokens, while reducing the row space along the direction of the embedding of $t^*$.
%
Conversely, when $t^*$ is being predicted ($t_i = t^*$), all token embeddings from the context are added to the row space of $\bm{W}_{qk}$, while only the embedding of $t^*$ is added to the column space.
%
Consequently, the role a token plays during training affects its contribution to the weight update differently.
%
We formalize this in the following proposition.
%
\input{ICML2025-rebuttal/math/statements/prop-gradient-columns-rows}

%
We provide proof for this proposition with related remarks in Appendix \ref{supp-math-prop-gradient-column-rows}.
%
 Note that all mathematical results derived so far rely solely on the structure of self-attention and make no assumptions about the input data.   
%
In the next section, we build on these results to link structural patterns in \( \bm{W}_{qk} \) with the specific form of the objective function.


\subsection{The relation between objective functions and structures in self-attention matrices}
%
In this final section, we show how to relate these properties to the specific objective function.
%
Crucially, the number of times a token \( t^* \) appears as context or as a prediction depends on the training objective.
%
Autoregressive training implicitly introduces directionality by predicting each token solely on its preceding tokens. 
%
In contrast, bidirectional training uses tokens as context and as predictions symmetrically.
%
This fundamental difference affects the weight updates of \( \bm{W}_{qk} \), and consequently, the structures encoded in its rows and columns.
%

We formalize this relationship in the next two theorems, under the following assumptions:  (a) tokens exhibit statistical correlations, leading to partial alignment in their embeddings; (b) the entries of \( \bm{W}_{qk} \) are i.i.d. at initialization with finite mean and variance; (c) some tokens tend to occur earlier in the sequence and are more predictive of future content; (d) bidirectional training induces approximately symmetric error signals.
%
In this setting, we prove how the objective function influences the internal structure encoded by self-attention.
%
\begin{theorem}
\label{theo-gradient-directionality}
%
(\textbf{Autoregressive training induces directionality})
Let \( \mathcal{V} = \{t_0, \dots, t_V\} \) be a vocabulary of tokens, and let \( \mathcal{U} \subset \mathcal{V}^N \) denote the sample space of all sequences of length \( N \). Let \( U = \{t_1, \dots, t_N\} \in \mathcal{U} \) be a random variable with \( U \sim P(U) \). Define \( \Pr[t_j = t^*] \) as the marginal probability that the token at position \( j \) equals \( t^* \in \mathcal{V} \).
%
Let \( \{\bm{x}_1, \dots, \bm{x}_N\} \) be the token embeddings corresponding to the elements of \( U \), where each embedding \( \bm{x}_i \sim \mathcal{D} \) is drawn i.i.d. from a distribution \( \mathcal{D} \) with zero mean and covariance matrix $\text{Cov}(\bm{x}_i)  = \Sigma$..
%
Let \( \bm{W}_{qk} \) be query-key matrix of a self-attention mechanism, and let \( \Delta \bm{W}_{qk} \) denote its gradient update as defined in Proposition~\ref{prop-gradients-self-attention}, computed under an autoregressive objective as in Definition~\ref{def-objective-functions}.
%

It follows that the contribution of the token \( t^* \) after training satisfies, 
%
\begin{equation}
\frac{\mathbb{E}_{\mathcal{D}}\left[ \|\Delta \bm{w}_{\cdot, k}\|^2 \right]}{\mathbb{E}_{\mathcal{D}}\left[ \|\Delta \bm{w}_{m, \cdot}\|^2 \right]} > 1 
\quad \forall k, \,\,\forall m\,\, \text{s.t.} \,\, \Sigma_{m,m} < \frac{\mathrm{Tr}(\Sigma)}{d} \,.
\end{equation}
%
Moreover, there exists a scalar \( \gamma \in \mathbb{R}_{>0} \), proportional to the product of the row and column norm variances, $
\gamma \propto \mathrm{Var}\left(\|\bm{w}_{m,\cdot}\|\right) \cdot \mathrm{Var}\left(\|\bm{w}_{\cdot,k}\|\right)$,
such that for all \( w > \gamma \),
\begin{equation}
    \Pr\left[\|\bm{w}_{\cdot, k}\|^2 > w\right] > \Pr\left[\|\bm{w}_{m, \cdot}\|^2 > w\right] \,,
\end{equation}
that is, columns of \( \bm{W}_{qk} \) are more likely than rows to have high norm under autoregressive training, indicating a directional bias in gradient updates.
%
\end{theorem}
%

\begin{theorem}
\label{theo-gradients-symmetry}
%
(\textbf{Bidirectional training induces symmetry}).
%
Let $U = \{t_1, \dots, t_N\}$ be a sequence of tokens.
%
Let \( \bm{W}_{qk} \) be the query-key matrix of a self-attention mechanism, and let \( \Delta \bm{W}_{qk} \) denote its gradient update as defined in Proposition~\ref{prop-gradients-self-attention}, computed under a bidirectional objective as in Definition~\ref{def-objective-functions}.
%
It follows that,
%
\begin{equation}
    \Delta \bm{W}_{qk} = \sum_{i = 1}^N \sum_{j = 1}^N \beta_{ij} \bm{K}_{ij} \,,
\end{equation}
%
and that every pair $(i,j)$ with $i \neq j$ contributes to the weight update with a term
%
\begin{equation}
    \Delta \bm{W}_{qk}\big|_{\bm{t}_i \leftrightarrow \bm{t}_j} = \beta_{ij}\bm{K}_{ij} + \beta_{ji}{\bm{K}_{ij}}^T \,,
\end{equation}
%
that is approximately symmetric,
%
\begin{equation}
    \Delta \bm{W}_{qk}\big|_{\bm{t}_i \leftrightarrow \bm{t}_j} \approx \Delta \bm{W}_{qk}\big|^\top_{\bm{t}_i \leftrightarrow \bm{t}_j} \,.
\end{equation}
%
\end{theorem}
%
We provide a proof of these two Theorems, and the related Propositions and Lemmas in Appendix \ref{supp-math-directionality-symmetry}.
%